\chapter{Conclusion and Recommendations}
\label{ch:conclusion}

\section{Introduction}
This chapter relates the implemented system and recorded evidence to the project objectives. It summarises achievements and difficulties, identifies justified future work, and gives a recommendation for the current local deployment.

\section{Conclusion}
The project produced a local visual fashion search application in which a user
can search a full photograph or communicate the intended garment through
interactive selection. SAM2.1 supplies point-guided cutouts, SigLIP2 supplies
fixed image representations, FAISS ranks normalised vectors, and FastAPI
connects those services to an inspectable browser workflow. An original
3,068-image research gallery and an independent managed-product catalogue
serve complementary experimental and practical roles.

The central finding is conditional rather than universal. The retained synthetic
test records a category Top-1 increase from 46\% to 75\% when the target garment
is selected, under its historical 100-query protocol. A different,
application-aligned clean-query comparison records a category Top-1 increase
from 92.20\% to 96.24\% on 372 test queries with automatic category filtering.
These results support the usefulness of target communication and category
control within their stated tasks. They do not establish exact-SKU accuracy or
general consumer-photo performance.

The seven-group identity pilot demonstrates retrieval of the reviewed alternate
views under explicit query exclusion, but its full-image and SAM2 reverse
conditions tie. Catalogue cropping lacks sufficient consistent benefit to
replace the original index. Keeping those conclusions distinct is part of the
project's contribution: a working application is accompanied by an account of
what its evidence does and does not support.
\subsection{Achievements}
The research objectives were addressed by comparing five pretrained encoders,
retaining separate selection and routing experiments, evaluating a candidate
gallery transformation, and reporting a reviewed identity pilot with explicit
limits. The deployed encoder was selected from the local comparison, while an
unsupported gallery change was not promoted.

The application objectives were addressed through full-image, SAM2, and manual
selection; accepted-query preview; selectable catalogue and category controls;
product CRUD with multiple photographs and availability; local administrator
protection; and backup, restore, and recovery. Final-report unit, JavaScript,
HTTP, administration, and evidence checks passed within their recorded scope.
The editable report connects these functions to requirements, architecture,
data structures, and results.

\subsection{Difficulties}
Several difficulties shaped the final design. A segmentation model may capture
only part of a garment, making multiple positive and negative points and manual
fallbacks necessary. Fine-grained visual resemblance is not fully represented
by a category label. Catalogue images can repeat products without providing
verified commercial identities. Comparing cached experiments with different
preprocessing can create apparently contradictory percentages unless their
protocols are separated.

Data administration introduced its own requirements: photograph updates must
remain consistent with vectors; stale editors must not replace newer content;
and restore must preserve a recoverable previous state. Runtime distribution
also requires checkpoints, images, an index, and dependencies in addition to
source code. A failed large-file transfer motivated locally verified split
packaging, while publication remains a separate delivery operation.

\subsection{Future Works}
\subsubsection{Verified Identities and Realistic Evaluation}
The first research priority is a catalogue with trustworthy product identities
and independently captured query views. Merchant-supplied SKU relationships or
a benchmark with defined cross-view identity annotations should be split by
the intended product/capture protocol before experimentation. Real multi-garment
customer-like images should accompany clean product photographs. Independent
judgments of shape, colour, and pattern resemblance can complement exact-identity
labels when alternative products are acceptable.

\subsubsection{Evidence-Based Ranking Improvements}
After establishing that evaluation set, compare the existing baseline against
candidate fine-grained reranking or domain adaptation. Any fusion of colour,
texture, global embeddings, or crop representations should be selected on
development data and tested once on held-out queries. The original index and
scoring path should remain available as a baseline. Similarity thresholds for
``no confident match'' require calibration against known present and absent
products rather than an arbitrary percentage.

\subsubsection{Selection Quality and User Experience}
A dedicated mask study should use independently annotated garment masks and
measure overlap and failure cases separately from retrieval. An evaluated
user study can measure task completion, refinement effort, and time to a useful
result. Further work may reduce repeated segmentation cost, improve small-item
selection, and evaluate the interface on multiple devices. These changes should
preserve the final-query preview and explicit acceptance that make errors visible.

\subsubsection{Store Deployment and Maintainability}
For a single store, the current managed catalogue provides a starting point.
Broader deployment would require server hardening, HTTPS, measured concurrency,
durable storage operations, and a consistent revision policy across all write
routes. A city-wide service additionally needs explicit store ownership,
seller roles, product-location/contact metadata, and controls for each store's
data. Payment and ordering would be separate requirements rather than assumed
extensions of image search.

The added catalogue currently constructs its exact index for each search.
If measured catalogue growth makes this costly, a versioned persistent index
or a suitable approximate search structure can be evaluated. Such a change
should preserve product-level aggregation and the separation of research and
managed data.

\section{Recommendation}
Use the current application as the documented local search baseline, complete
its runtime delivery when the connection permits, and prioritise trustworthy
evaluation data before adding more ranking complexity. Retain the original
gallery index until a controlled comparison supports a replacement. Subsequent
changes can then be assessed against a clear software and research reference.
